\section{FINDINGS}
\label{sec:findings}


\begin{figure*}[t]
  \centering
  \includegraphics[width=0.243\linewidth]{Figures/runs-narrowroad.jpg}\hfill
  \includegraphics[width=0.243\linewidth]{Figures/runs-corner.jpg}\hfill
  \includegraphics[width=0.243\linewidth]{Figures/runs-crossroad.jpg}\hfill
  \includegraphics[width=0.243\linewidth]{Figures/runs-outofstore.jpg}
  \caption{All P1--P13 robot runs drawn over each scene's top-down capture, showing from left to right the narrow road, corner interaction, crossroad, and out of store. Blue is the autonomous planner, orange is the participant-driven runs, $\times$ is the robot goal, and each scale bar is 5\,m. The planner repeats essentially one route per scene while participant-driven routes fan out around it. They pull to the sidewalk edges on the narrow road, swing wide of the blind corner, and cross straight before turning at the crossroad, the same legibility norms participants stated in the interview.}
  \Description{Four overhead scene photographs with robot trajectories overlaid, one per scenario. In each, the autonomous planner's blue lines form a single tight bundle, while orange participant-driven lines fan out around it. On the narrow road the orange routes hug the sidewalk edges where the planner holds the centre. At the corner they arc wide of the planner's tight inside line; at the crossroad they cross straight along the crossing band and turn afterwards where the planner cuts a diagonal. Out of the store they spread around the doorway area the planner drives through. Each panel carries a five-metre scale bar and a black cross at the robot goal.}
  \label{fig:runs}
\end{figure*}
\begin{table}[t]
  \centering
  \small
  \setlength{\tabcolsep}{4pt}
  \caption{Autonomous planner vs.\ participant driving}
  \label{tab:ladder}
  \begin{tabular}{lrrrr}
    \toprule
     & \textbf{Stops $\geq$1\,s} & \textbf{Stopped} & \textbf{Speed} & \textbf{Path}\\
     & \textbf{near human} & \textbf{(s)} & \textbf{near/free} & \textbf{eff.}\\
    \midrule
    Auto policy (DWA)   & 0.00 & 0.2 & 0.93 & 0.99\\
    User demonstration  & 0.67 & 4.5 & 0.83 & 0.94\\
    \bottomrule
  \end{tabular}

  \vspace{2pt}
  \parbox{0.92\columnwidth}{\footnotesize Speed is the mean speed near a human over the mean free speed, and efficiency is the median. Trailing idle at the goal is excluded from stop measures.}
\end{table}
\subsection{How do participants perceive social navigation experiences and identify disruption?}
The two largest variances in the rating sheet are the blind corner on comfort ($\sigma^2{=}7.81$) and the narrow road on path adjustment ($\sigma^2{=}7.21$). Some participants had to adjust their paths to accommodate the robot, and that adjustment is what degraded their experience. The trajectories show the same disruption: on the narrow road, where the participant meets the robot head-on, the median participant drops to $0.40\times$ their own free-walking speed at closest approach and stops 4.4 times per trial.

The go-along replay traces the reason for the disruption. The most prevalent one is being pushed onto ground they would never normally walk on. P6 stepped onto a hatchway \emph{on purpose} to escape the robot, scoring comfort 2 with adjustment 10 (\emph{``it was coming too fast that I couldn't go around''}). P9 named the trade-off (\emph{``if I need to get to my destination now, I'm more likely to step on something I wouldn't normally step on''}), and, asked whether ceding the path frustrated her, answered \emph{``Probably. Only because it's not human.''}

Participants who trusted the robot's sensing walked on and were fine: they did not adjust their path but assumed the robot would yield. P12 proceeded because \emph{``this kind of robot has a camera,''} then admitted she watched it anyway (\emph{``I was a little worried about whether the robot successfully detected me''}). Comfort and disruption also come apart in the ratings. P7 adjusted \emph{``just as much as I would another person''} and rated comfort 10 (\emph{``that's how people walk too''}). P5, in contrast, adjusted her path \emph{zero} at the crossroad yet rated comfort 4, evidence of a vigilance tax with no detour attached (\emph{``it's still in your line of sight, so you also have to be cautious of it''}).

Whichever prior they held, participants shared one underlying behavior: they pre-emptively cleared out of the robot's way because they did not trust it to react. P4 put it directly, \emph{``I wanted to avoid the robot, to get out of its path, as opposed to having it recognize me and move.''} Most of the disturbance came not from yielding, which participants grant other people all day, but from the uncertainty underneath it (P8, \emph{``I can't predict where it's actually gonna go''}) and from unreadable capability (P13, a person is \emph{``probably as agile as I am\dots\ with the robot, I'm not sure how smart it is''}). Participants also borrow a frame for the robot, whether a vehicle (P1's \emph{``mini car''}), a pet (P9's \emph{``dog off the leash\dots\ better, because it can't fight me''}), or another pedestrian (P7, P12). The frame predicts the ruleset. P7, holding the pedestrian frame, rated every scenario no worse than a human doing the same thing, and the vehicle framers imported traffic norms wholesale.

Underneath these reactions sits the actual navigational challenge: reading the other agent's next move from limited information, and it binds both directions of the encounter. Driving the robot, P5 took the efficient shortcut and recanted it within a minute, and P13 steered around a stopped pedestrian, then judged herself on the replay (\emph{``now thinking back, it was a little bit rude''}). P9 switched into the robot's first-person camera and found the perception problem harder than she had assumed (\emph{``this is much scarier\dots\ looking at people's knees''}), driving differently for the rest of her session. The world-editing tool isolates the variable. When the pedestrian was forced onto a known straight path, P3's whole problem dissolved (\emph{``way more relaxed\dots\ you do what you gotta do, and I'll do what I gotta do''}), while a faster, diagonal crossing collapsed P11's comfort from 8 to 3--4 (\emph{``it squeezed me off the edge --- why is this robot trying to take the road from me?''}). Editing the same scene in both directions separates the variable: the challenge is not the proximity of the robot but the predictability of its behavior.

\subsection{What are some social navigation strategies and considerations that we can learn?}
Participants reason about the robot scenario by scenario, and both demonstration channels turn that reasoning into visible trajectories: the joystick in the moment and the redraw tool on reflection. At the blind corner, most participants recognized a trade-off between efficiency and mutual visibility, and visibility won. P5 took the efficient inside line and then recanted unprompted (\emph{``socially it would have been more acceptable to swing out wider''}). P8 drove the wide arc deliberately (\emph{``larger curvature so the human can see me''}). P10 rejected the inside line even where keep-right favored it (\emph{``you're not gonna see a person until the last minute''}). P11 split the turn into two straight legs with a pivot in place, so every segment stays extrapolatable. At the head-on encounter, the reasoning shifts to stop placement, because where you stop is the message. Participants stopped out of the flow but still facing the person, on stable ground rather than the grate, and out of the door-swing arc (P9, \emph{``no point being out of the way, only to be getting hit by doors''}). P13 wrapped it all in a motion policy, \emph{``go straight, but go slowly\dots\ set the right expectation that I won't move wildly.''}

Table~\ref{tab:ladder} contrasts the two. SEAN's default DWA planner produces no stop of at least one second near a human across all 74 autonomous runs. However, \textbf{\textit{participants stop routinely and treat stopping as a preferred behavior for social navigation.}} The routes participants \emph{drew} ($n{=}36$) push every measure further still, with 0.83 stops near a human, 7.2\,s stopped, a 0.77 speed ratio, and 0.90 efficiency. Driven paths agree with the planner only where nothing social happens (mean deviation 0.39\,m, Fr\'{e}chet distance 1.19\,m at the median trial, and 67\% of a typical driven path within half a metre of the planner's line), and the disagreement concentrates at encounters (Fig.~\ref{fig:runs}).

Participants do not converge on one another either. Routes that different participants drove in the same scene share a median corridor IoU of only 0.56 (0.52--0.59 across scenes), so for demonstration-based learning the human target is a \emph{distribution} over legible routes, not a single trajectory to imitate.

Three social norms recur across participants. \emph{Follow people}: in the crossroad scenario a walker is a path-clearer, an oracle, and a shield (P8, \emph{``the human makes smarter decisions --- trail it, copy what it does''}), and P11 applied it at the crossing (\emph{``if the person goes, I follow them --- the driver can see a person and won't dare hit them''}). \emph{Do not match human walking speed}: as P7 put it, \emph{``it should either be so fast that I don't have to wait for it, or so slow that I can bypass it.''} \emph{Wait, but with a cap, then escalate}: P13 imposed the bound on herself mid-drive (\emph{``there should probably be a cap waiting time''}), and P8 supplied the escalation, \emph{``wait a couple seconds, signal; if nothing happens, take the other route.''} Strategies also changed live inside the simulation, the concrete benefit of an interactive instrument over an interview. P3's yielding policy rewrote itself the moment he learned the hatch could trap his wheels (\emph{``I could theoretically get stuck in that\dots\ then I'd wait for the person to pass me''}). P8 converted the uncertainty about terrain directly into a reason to yield. P5, boxed in by edited clutter, learned passability by watching the simulated pedestrian thread the cones and copied the route. P10, after blocking the wide side of his own corner, produced the study's clearest courtesy gesture --- a step back to cede the lane (\emph{``a gesture that says: go first''}).

% \subsection{Signaling: convergent structure, contested register}
% Signaling proposals converged on a structure rather than a modality. Participants routed the channel by geometry and noise, visual when in the person's field of view and sound when behind them (P8), voice where it is quiet and light where it is loud (P12, \emph{``how we predict a car is to see the lights''}). They gated announcements by proximity, density, or a failed wait (P12, \emph{``if the human is within a certain distance, then say something --- not all the time''}). And they asked that the yield itself be legible. P12's \emph{``I'm waiting for you''} doubles as proof of detection (\emph{``then it's super clear I can just go''}). P10 demands a stop that does not read as a breakdown (\emph{``not done, taking a break, or its battery is dead''}). P13 adds sequencing (\emph{``say excuse me, then wait one or two seconds, for the pedestrian to be aware''}).

% The register, however, is contested, and the splits read as personalization axes rather than design defects. \emph{``Excuse me''} is rude presumption to P2 (\emph{``it doesn't mean I want to excuse you''}) and P4, yet perfectly welcome to P12 and P13. Voice is \emph{``creepy''} to P2 and \emph{``weird''} to P10, yet delight on the receiving end when role-played (P12, \emph{``oh! The robot is talking!''}), and P4 offered the synthesis of \emph{``a robot voice, but a kind robot voice.''} Whether the robot should travel in front or behind splits the same way. P4 and P11 want it behind (\emph{``I can just forget about it''}); P6, P8, and P9 want it visible in front (\emph{``I don't trust it enough when I can't see it''}). Every signal design also presumes three things that all break at scale. It presumes compliance (P8, \emph{``indicators would work, but if people don't adhere to it, what's the point?''}), attribution (P10, with multiple robots \emph{``humans can't distinguish which one is giving the signal''}), and sensory ability, since deaf pedestrians need the visual channel and blind pedestrians the audio one, and P9 notes service animals are not trained on robots at all.

\subsection{What are some challenging scenarios that they came up with?}

\begin{figure}[t]
  \centering
  \includegraphics[width=\columnwidth]{Figures/wb-examples.jpg}
  \caption{Examples from the worlds participants built, clockwise from top left: a cyclist, a wheelchair user, a table, and a parked scooter added to the narrow road, with the robot's goal lettered on the storefront; a top-down view of the corner, where added planters, bins, and litter pinch the walkway and the drawn robot and pedestrian paths thread the remaining gap; the crossroad, where a car takes the crossing as the robot waits at the stop sign with its goal across the street; and the store exit, where a crowd fills the sidewalk between the robot and the marked pedestrian goal. Most of what participants placed is a person or a wheeled user rather than an obstacle.}
  \Description{Four simulator screenshots in a two-by-two grid. Top left: an elevated view of a narrow sidewalk where a woman cycles along the curb, a man in a wheelchair sits mid-sidewalk inside a green selection box, and a table, a parked scooter, a child, and scattered litter fill the walkway, with gold Robot Goal lettering on the wall. Top right: a top-down view of a street corner where planters, trash bins, and a fire hydrant line the building edge, and a red and a magenta path line weave between them toward a doorway past a crosswalk. Bottom right: a street-level view from behind an avatar at a crosswalk, with a STOP sign at center, a silver car entering the crossing from the left, and a gold Robot Goal label on the far corner. Bottom left: an elevated view over a store entrance where pedestrians cluster around the orange door, a cyclist passes on the street, drawn red and magenta paths cross the pavement, and a black Pedestrian Goal banner spans the sidewalk.}
  \label{fig:wb-examples}
\end{figure}

% \begin{figure*}[t]
%   \centering
%   \includegraphics[width=\linewidth]{Figures/wb-overlap.jpg}
%   \caption{Where P1--P13 edited the world, over each scene's top-down capture (grey). Heat counts the distinct participants who placed or moved an object within 2\,m of a spot (latest save per participant and scene); dots are individual placements (blue: human agent; orange: obstacle or prop; green: ground decal), and open squares are pre-existing objects moved over 1\,m. Without seeing one another's sessions, 10 of the narrow road's 12 builders converge on the same bottleneck and 10 of the corner's 11 on the same approach, while the crossroad --- edited through geometry and behavior rather than clutter --- stays sparse.}
%   \Description{Four grey overhead scene photographs in a row, one per scenario, each overlaid with a blue heatmap and colored dots marking placed objects. On the narrow road and the corner the heat concentrates into a single dark region annotated ten of twelve and ten of eleven participants respectively; the crossroad shows several small light patches annotated five of five; the out-of-store panel shows one concentrated region annotated seven of eight. A shared color scale from one to ten participants and a marker legend sit below the panels, and each panel carries a five-metre scale bar.}
%   \label{fig:worlds}
% \end{figure*}

Asked to make each scene harder, participants placed 143 objects across 39 edited worlds (Fig.~\ref{fig:wb-examples}). 58\% of everything placed is a person with a specific identity --- pedestrians, cyclists, wheelchair users --- so they build social pressure, not obstacle courses. Sixteen mobility-aid users (wheelchairs, mobility scooters, and cane users) were placed unprompted by roughly half the panel, indicating a general concern that sidewalk robots must share the walkway with people whose mobility, speed, and room to maneuver differ from the able-bodied pedestrian most planners assume.

Placement also converged spatially, even though no participant saw another's session. At the hottest spot of the narrow road, 10 of its 12 builders placed or moved something within 2\,m, and 10 of 11 did at the corner. The corner edits share one intent: blocking the sightline to the other sidewalk, so that robot and pedestrian meet with no warning.

Another common pattern is blocking the road itself. In the head-on encounter, P2 narrated the design as she built it, \emph{``put something right opposite the table, where there's only space for one person, and we get there at the same time --- like, what are we gonna do?''} At the crossroad participants ideated on behavior instead --- what if the car starts to turn right into the crossing --- and the researcher Wizard-of-Oz'd the car live (Fig.~\ref{fig:wb-examples}, bottom right). Participants also staged specific social dilemmas. P11 reconstructed scaffolding whose pole rows \emph{``the robot can pass under but a person can't.''} P10 recreated a reported case of a robot camped on a curb cut while a wheelchair user's crossing timer ran out. P9 placed cones because they get kicked (\emph{``how do you get a robot to navigate around something that's in the wrong spot?''}), and P6 requested the study's only weather edit --- rain, \emph{``because it can be hard to see the robot.''}

\smallskip\noindent\textbf{Summary.}
The findings form one arc. Participants' discomfort traces to the robot's predictability rather than its proximity, and to the ground they surrender when they do not trust it to react. Their demonstrations translate that demand into concrete, learnable geometry --- wide visible arcs, meaningful stop placement, speeds offset from human walking pace --- that the default planner never produced, and they form a distribution over legible routes rather than one trajectory. Handed the world itself, participants spent their effort manufacturing social pressure and dilemmas, from blocked sightlines to contested crossings and curb cuts, rather than physical difficulty. The simulator's contribution is that all three surfaced as observable behavior, rated, driven, drawn, and built, instead of merely stated.

% Two results bear directly on the simulator. First, 83\% of the participants placed static agents, and static agents were immediately reclassified as furniture (P4 \emph{``treated them as regular obstacles to go around, instead of having behavior to consider''}), so social world building needs moving agents by default. Second, re-behaving existing agents proved as analytically productive as placing new objects. The non-yielding pedestrian forced P11's robot into a deadlock. P7 read the erratic mirror-walker instantly as \emph{``how an old person would react,''} and speed swaps ran in both directions. These edits form paired manipulations, such as erratic against predictable and yielding against non-yielding, which license separating predictability from proximity above. Behavior-level editing therefore deserves first-class interface support rather than wizard improvisation.

% \subsection{Deployment concerns beyond any single policy}
% Once participants stopped driving and started imagining real streets, three concerns recurred that no navigation policy addresses alone. The first is \emph{cars}. Eight participants independently raised that drivers cannot see something knee-high (P8, at the crosswalk as the robot, \emph{``I wasn't sure if the car would have seen me\dots\ I don't know if the car's gonna see the small robot''}), and the shared mitigation is to cross behind a human the driver can see. The second is \emph{harassment, theft, and the urgency-disclosure dilemma}, reconstructed unprompted by four participants (P2, P3, P4, P8). Negotiating priority requires broadcasting the mission, yet the broadcast invites theft (P4, \emph{``maybe if it was carrying medicine it could say so --- but then what if people try to steal it, if they know it's inside?''}), and the dilemma remains unresolved by every participant who found it. P4's field observation cuts the other way, since harassment is curiosity that starves without a reaction (\emph{``they lose interest pretty quickly if a robot doesn't do anything''}). The third is \emph{density ceilings}. P11 had seen the end state (\emph{``in Berkeley I once saw three or four robots in a row, all stuck on a long narrow sidewalk''}), and P7 supplied the geographic irony that robots are most useful exactly where density defeats them (\emph{``in New York it's gonna be helpful, because being a delivery driver in New York is a nightmare''}). Half the panel reached for a dedicated robot lane as the exit, and the other half dismantled it as premature, a theft of bus-and-bike space, and a relocation of the conflict to every crossing. Participants want the separation, and nobody can find the space.

% The cleanest theoretical frame came from the three participants who had watched service robots succeed elsewhere. Restaurant robots work, P12 argued, because \emph{``it's a confined space, and people basically have one intention, which is to eat --- public space is more mysterious.''} P8 supplied the mechanism, observing that the robot \emph{``was pretty predictable --- and it only worked if things around it and the environment also worked in a predictable way.''} The sidewalk is hard because it is a maximum-entropy intention environment --- which is also why the predictability finding above, rather than any single behavior, is the design target.
